\section*{Bab \ref{ch:g6:rocks-and-minerals} --- Batuan dan Mineral}

\begin{solution}{exo:g6:rocks-and-minerals:1}
Granit: beku. Batu pasir: sedimen. Marmer: metamorf. Basal: beku. Batu
kapur: sedimen. Batu sabak: metamorf.
\end{solution}

\begin{solution}{exo:g6:rocks-and-minerals:2}
Kuarsa adalah \omterm{def:g6:rocks-and-minerals:mineral}{mineral}; granit adalah \omterm{def:g3:stones-pebbles-sand:rock}{batuan}. \omterm{def:g6:rocks-and-minerals:mineral}{Mineral} adalah satu zat alami
yang atom-atomnya tersusun secara teratur; \omterm{def:g3:stones-pebbles-sand:rock}{batuan} tersusun dari satu atau
beberapa \omterm{def:g6:rocks-and-minerals:mineral}{mineral} yang melekat menjadi satu (granit: feldspar, kuarsa, dan
mika).
\end{solution}

\begin{solution}{exo:g6:rocks-and-minerals:3}
Talk (1), kalsit (3), kuarsa (7), intan (10). Intan, yang paling keras,
menggores semua \omterm{def:g6:rocks-and-minerals:mineral}{mineral} lainnya.
\end{solution}

\begin{solution}{exo:g6:rocks-and-minerals:4}
``\omterm{def:g3:stones-pebbles-sand:rock}{Batuan} yang meleleh menjadi \omterm{def:g6:rocks-and-minerals:rock-families}{batuan beku} melalui pendinginan. Batu kapur
menjadi marmer melalui panas dan tekanan. Granit berubah menjadi butiran
\omterm{def:g3:stones-pebbles-sand:sand}{pasir} melalui \omterm{def:g5:water-shapes-the-land:weathering}{pelapukan}.''
\end{solution}

\begin{solution}{exo:g6:rocks-and-minerals:5}
(a) Basal. (b) Marmer. (c) Batu pasir.
\end{solution}

\begin{solution}{exo:g6:rocks-and-minerals:6}
\omterm{def:g6:rocks-and-minerals:hardness}{Kekerasannya} lebih besar daripada \omterm{def:g6:rocks-and-minerals:hardness}{kekerasan} kuku jari (kira-kira 2.5) dan
lebih kecil daripada \omterm{def:g6:rocks-and-minerals:hardness}{kekerasan} koin tembaga (kira-kira 3.5). \omterm{def:g6:rocks-and-minerals:mineral}{Mineral} itu
mungkin kalsit (3).
\end{solution}

\begin{solution}{exo:g6:rocks-and-minerals:7}
Kuarsa: $21 \div 60 = 0.35$, yaitu 35~\%. Feldspar: $33 \div 60 = 0.55$,
yaitu 55~\%. Mika: $6 \div 60 = 0.10$, yaitu 10~\%. Periksa:
$35 + 55 + 10 = 100$.
\end{solution}

\begin{solution}{exo:g6:rocks-and-minerals:8}
Granit mendingin dengan sangat lambat, jauh di bawah tanah: butirannya
mempunyai waktu yang lama untuk tumbuh, dan menjadi besar. Basal berasal
dari \omterm{def:g4:volcanoes-earthquakes-first-look:eruption}{lava}, yang mendingin dengan cepat di permukaan: butirannya tidak
mempunyai waktu untuk tumbuh dan tetap terlalu kecil untuk dilihat.
\end{solution}

\begin{solution}{exo:g6:rocks-and-minerals:9}
(a) Batu kapur adalah \omterm{def:g6:rocks-and-minerals:rock-families}{batuan sedimen}; marmer adalah \omterm{def:g6:rocks-and-minerals:rock-families}{batuan metamorf}.
(b) Jauh di bawah tanah, panas dan tekanan mengubahnya tanpa melelehkannya:
kalsitnya tumbuh menjadi kristal-kristal mengilap yang saling mengunci.
(c) Keduanya tersusun dari \omterm{def:g6:rocks-and-minerals:mineral}{mineral} yang sama, yaitu kalsit, dan kalsitlah
yang berbuih jika ditetesi asam.
\end{solution}

\begin{solution}{exo:g6:rocks-and-minerals:10}
(a) Ya: kuarsa (7) lebih keras daripada kaca (kira-kira 5.5). (b) Tidak:
fluorit (4) lebih lunak daripada kaca. (c) Ya: paku baja (kira-kira 6.5)
lebih keras daripada feldspar (6).
\end{solution}

\begin{solution}{exo:g6:rocks-and-minerals:11}
Misalnya: (1) granit \omterm{def:g1:landscapes-around-us:landform}{gunung} itu (\omterm{def:g6:rocks-and-minerals:rock-families}{batuan beku}) dihancurkan oleh \omterm{def:g5:water-shapes-the-land:weathering}{pelapukan},
dan butiran kuarsanya terlepas; (2) \omterm{def:g1:landscapes-around-us:water}{sungai} mengangkutnya pergi: \omterm{def:g5:water-shapes-the-land:erosion}{erosi};
(3) butiran itu diendapkan berlapis-lapis di dasar \omterm{def:g2:water-on-earth:ocean}{laut}: \omterm{def:g5:water-shapes-the-land:deposition}{pengendapan};
(4) setelah terkubur di bawah lapisan-lapisan lain, butiran itu tertekan
dan tersemen menjadi batu pasir: pengerasan (\omterm{def:g3:stones-pebbles-sand:rock}{batuannya} sekarang \omterm{def:g6:rocks-and-minerals:rock-families}{batuan
sedimen}); (5) setelah terbawa jauh ke bawah tanah, batu pasir itu diubah
oleh panas dan tekanan menjadi \omterm{def:g6:rocks-and-minerals:rock-families}{batuan metamorf}.
\end{solution}

\begin{solution}{exo:g6:rocks-and-minerals:12}
(a) Tidak tergores oleh baja (kira-kira 6.5) dan dapat menggores kaca:
lebih keras daripada 6.5. Tergores oleh topaz (8): lebih lunak daripada 8.
\omterm{def:g6:rocks-and-minerals:hardness}{Kekerasannya} 7: \omterm{def:g6:rocks-and-minerals:mineral}{mineral} itu kuarsa. (b) Kuarsa lebih keras daripada hampir
segala sesuatu yang bergesekan dengannya (butiran lain, \omterm{def:g3:stones-pebbles-sand:pebble}{kerakal}): kuarsa
hampir tidak aus. Kalsit (3) dan mika (2 sampai 3) jauh lebih lunak:
gesekan menggerusnya menjadi kepingan kecil yang terbawa pergi.
\end{solution}

\begin{solution}{pb:g6:rocks-and-minerals:1}
\textbf{1.} Granit, \omterm{def:g6:rocks-and-minerals:rock-families}{batuan beku}; \omterm{def:g6:rocks-and-minerals:mineral}{mineralnya} adalah feldspar (merah muda atau
putih), kuarsa (kelabu, seperti kaca), dan mika (hitam, mengilap).

\textbf{2.} Batu kapur, \omterm{def:g6:rocks-and-minerals:rock-families}{batuan sedimen}. Cangkang-cangkang itu dulu milik
hewan yang hidup di \omterm{def:g2:water-on-earth:ocean}{laut}; sisa-sisanya diendapkan dan mengeras menjadi
\omterm{def:g3:stones-pebbles-sand:rock}{batuan}.

\textbf{3.} Batu sabak, \omterm{def:g6:rocks-and-minerals:rock-families}{batuan metamorf}, yang terbentuk dari serpih.

\textbf{4.} Marmer, \omterm{def:g6:rocks-and-minerals:rock-families}{batuan metamorf}, yang terbentuk dari batu kapur.

\textbf{5.} Kuarsa: tidak tergores oleh baja (kira-kira 6.5), menggores
kaca (kira-kira 5.5); \omterm{def:g6:rocks-and-minerals:hardness}{kekerasannya} 7.

\textbf{6.} Tergores oleh baja (kira-kira 6.5) tetapi tidak oleh kaca
(kira-kira 5.5): antara 5.5 dan 6.5. \omterm{def:g6:rocks-and-minerals:mineral}{Mineral} itu feldspar (ortoklas, 6).

\textbf{7.} Mika (\omterm{def:g6:rocks-and-minerals:hardness}{kekerasan} 2 sampai 3).

\textbf{8.} Ya: ubin jalan itu batu kapur, yang tersusun dari kalsit (3),
lebih lunak daripada koin tembaga (kira-kira 3.5).

\textbf{9.} \omterm{def:g5:water-shapes-the-land:weathering}{Pelapukan}: \omterm{def:g4:clouds-rain-snow:precipitation}{hujan} perlahan-lahan mengikis permukaan batu itu.

\textbf{10.} \omterm{def:g1:landscapes-around-us:landform}{Gunung} itu melapuk; butiran kuarsanya diangkut oleh \omterm{def:g1:landscapes-around-us:water}{sungai} dan
diendapkan berlapis-lapis, di dasar \omterm{def:g1:landscapes-around-us:water}{danau} atau \omterm{def:g2:water-on-earth:ocean}{laut}; setelah terkubur dan
mengeras, butiran itu dapat menjadi batu pasir, sebuah \omterm{def:g6:rocks-and-minerals:rock-families}{batuan sedimen}.

\textbf{11.} Ketiga keluarga ada: \omterm{def:g6:rocks-and-minerals:rock-families}{batuan beku} (granit air mancur), \omterm{def:g6:rocks-and-minerals:rock-families}{batuan
sedimen} (batu kapur ubin jalan), \omterm{def:g6:rocks-and-minerals:rock-families}{batuan metamorf} (batu sabak atap dan
marmer patung).

\textbf{12.} \omterm{def:g6:rocks-and-minerals:mineral}{Mineral} yang paling keras adalah kuarsa dalam granit:
\omterm{def:g6:rocks-and-minerals:hardness}{kekerasan} Mohs \textbf{7}.
\end{solution}
